
\documentclass[11pt]{article}
\usepackage[a4paper,margin=3cm]{geometry}
\usepackage[T1]{fontenc}
\usepackage[utf8]{inputenc}
\usepackage[spanish,es-nodecimaldot]{babel}
\usepackage{lmodern,microtype}
\usepackage{amsmath,amssymb,amsthm,mathtools}
\usepackage{booktabs,array}
\usepackage{hyperref}
\usepackage{fancyhdr}
\hypersetup{colorlinks=true,linkcolor=black,urlcolor=blue}
\setlength{\parindent}{0pt}\setlength{\parskip}{5pt}
\newtheorem{teorema}{Teorema}[section]
\newtheorem{proposicion}[teorema]{Proposición}
\theoremstyle{definition}
\newtheorem{definicion}[teorema]{Definición}
\newtheorem{ejemplo}[teorema]{Ejemplo}
\newcommand{\R}{\mathbb R}
\newcommand{\D}{\frac{d}{dx}}
\pagestyle{fancy}\fancyhf{}\lhead{\small Ecuaciones diferenciales de orden superior}\rhead{\small\thepage}
\title{\textsc{Ecuaciones Diferenciales I}\\[4pt]\Large Ecuaciones diferenciales de orden superior\\[5pt]\large Una guía de ejemplos y métodos}
\author{}\date{}
\begin{document}
\maketitle
\begin{center}\emph{Métodos, ejemplos resueltos y práctica.}\end{center}

\section*{Introducción}
Una ecuación diferencial de orden superior relaciona una función desconocida con sus
derivadas hasta un orden $n\geq2$. En la práctica, la pregunta fundamental es:
\emph{¿qué estructura tiene la ecuación y qué método permite encontrar sus soluciones?}

Estas notas están organizadas alrededor de esa pregunta. Se introducen las herramientas
mínimas y, de inmediato, se aplican a ejemplos de dificultad creciente. La teoría aparece
sólo en la medida necesaria para entender por qué funciona cada procedimiento.

\tableofcontents

\section{La forma general}
Una ecuación de orden $n$ puede escribirse como
\[
F(x,y,y',\ldots,y^{(n)})=0.
\]
En el caso lineal:
\[
a_n(x)y^{(n)}+\cdots+a_1(x)y'+a_0(x)y=g(x).
\]
Si $g=0$ es homogénea; si $g\ne0$, no homogénea.

\begin{ejemplo}
Clasifica:
\[
y''-3y'+2y=0,\qquad x^2y''+xy'-4y=x^3,\qquad y''+\sin y=0.
\]
Las dos primeras son lineales; la primera es homogénea, la segunda no homogénea.
La tercera es no lineal.
\end{ejemplo}

\section{Problemas de valores iniciales}
Para orden $n$ se prescriben normalmente
\[
y(x_0)=y_0,\quad y'(x_0)=y_1,\quad\ldots,\quad y^{(n-1)}(x_0)=y_{n-1}.
\]

\begin{teorema}[Existencia y unicidad]
Si los coeficientes de una ecuación lineal son continuos en un intervalo $I$ y el coeficiente principal no se anula en $I$, el problema de valores iniciales tiene una única solución en $I$ para datos iniciales dados.
\end{teorema}

\begin{ejemplo}[PVI completo]
Resolver
\[
y''-4y=12x,\qquad y(0)=4,\quad y'(0)=1.
\]
La homogénea da $y_h=c_1e^{2x}+c_2e^{-2x}$. Probamos $y_p=Ax+B$ y obtenemos
$A=-3$, $B=0$. Por tanto
\[
y=c_1e^{2x}+c_2e^{-2x}-3x.
\]
Las condiciones iniciales producen
\[
c_1+c_2=4,\qquad 2c_1-2c_2-3=1,
\]
de donde $c_1=3$, $c_2=1$. Así,
\[
\boxed{y=3e^{2x}+e^{-2x}-3x}.
\]
\end{ejemplo}

\begin{ejemplo}[Oscilador armónico]
Resolver
\[
x''+\omega^2x=0,\qquad x(0)=x_0,\quad x'(0)=v_0.
\]
Como $x=C_1\cos(\omega t)+C_2\sin(\omega t)$,
\[
\boxed{x(t)=x_0\cos(\omega t)+\frac{v_0}{\omega}\sin(\omega t)}.
\]
\end{ejemplo}

\section{Problemas de valores en la frontera}
Se prescriben condiciones en puntos distintos, por ejemplo
\[
y(a)=\alpha,\qquad y(b)=\beta.
\]
Puede haber cero, una o infinitas soluciones.

\begin{ejemplo}[Las tres posibilidades]
Sea
\[
y=c_1x^2+c_2x^4+3.
\]
Para $y(-1)=0$, $y(1)=4$ se obtienen $c_1+c_2=-3$ y $c_1+c_2=1$: no existe solución.

Para $y(0)=3$, $y(1)=0$ sólo se exige $c_1+c_2=-3$: hay infinitas soluciones.

Para $y(1)=3$, $y(2)=15$:
\[
c_1+c_2=0,\qquad 4c_1+16c_2=12,
\]
de donde $c_1=-1$, $c_2=1$ y
\[
\boxed{y=x^4-x^2+3}.
\]
\end{ejemplo}

\section{El operador diferencial}
Definimos
\[
D=\frac{d}{dx}.
\]
Entonces $D^2y=y''$, $D^3y=y'''$, etc. Por ejemplo,
\[
y''-2y'+5y=0
\quad\Longleftrightarrow\quad
(D^2-2D+5)y=0.
\]

\begin{ejemplo}[Factorización]
\[
y''-5y'+6y=0
\quad\Longleftrightarrow\quad
(D-2)(D-3)y=0.
\]
Por tanto
\[
\boxed{y=c_1e^{2x}+c_2e^{3x}}.
\]
\end{ejemplo}

\section{Superposición}
\begin{teorema}[Principio de superposición]
Si $y_1,\ldots,y_k$ son soluciones de una ecuación lineal homogénea, entonces
\[
y=c_1y_1+\cdots+c_ky_k
\]
también es solución.
\end{teorema}

\begin{ejemplo}
Para
\[
y''+2y'+y=0
\]
son soluciones $e^{-x}$ y $xe^{-x}$. Luego
\[
\boxed{y=c_1e^{-x}+c_2xe^{-x}}.
\]
\end{ejemplo}

\section{Independencia lineal y Wronskiano}
Las funciones $f_1,\ldots,f_n$ son linealmente independientes si
\[
c_1f_1+\cdots+c_nf_n=0
\]
implica $c_1=\cdots=c_n=0$.

Para dos funciones:
\[
W(f_1,f_2)=
\begin{vmatrix}f_1&f_2\\f_1'&f_2'\end{vmatrix}
=f_1f_2'-f_1'f_2.
\]
Si $W(x_0)\ne0$, son independientes.

\begin{ejemplo}
Para $f_1=e^{2x}$ y $f_2=e^{-x}$,
\[
W=-3e^x\ne0.
\]
Luego son linealmente independientes en $\R$.
\end{ejemplo}

\begin{ejemplo}
\[
f_1=4x^3,\quad f_2=2x^2,\quad f_3=8x^3+12x^2.
\]
Como
\[
f_3=2f_1+6f_2,
\]
el conjunto es dependiente.
\end{ejemplo}

\section{Conjuntos fundamentales}
Un conjunto de $n$ soluciones linealmente independientes de una ecuación homogénea de orden $n$ es un conjunto fundamental. Si $\{y_1,\ldots,y_n\}$ es fundamental,
\[
\boxed{y=c_1y_1+\cdots+c_ny_n}.
\]

\section{Coeficientes constantes}
Para
\[
a_ny^{(n)}+\cdots+a_1y'+a_0y=0
\]
probamos $y=e^{rx}$. Aparece el polinomio característico
\[
p(r)=a_nr^n+\cdots+a_1r+a_0.
\]

\subsection{Raíces reales distintas}
\begin{ejemplo}
\[
y''-5y'+6y=0.
\]
Como $(r-2)(r-3)=0$,
\[
\boxed{y=c_1e^{2x}+c_2e^{3x}}.
\]
\end{ejemplo}

\begin{ejemplo}
\[
y'''-y''-6y'+6y=0.
\]
El polinomio es $(r-1)(r^2-6)$, así que
\[
\boxed{y=c_1e^x+c_2e^{\sqrt6x}+c_3e^{-\sqrt6x}}.
\]
\end{ejemplo}

\subsection{Raíces repetidas}
Si $r$ tiene multiplicidad $m$, aparecen
\[
e^{rx},xe^{rx},\ldots,x^{m-1}e^{rx}.
\]

\begin{ejemplo}
\[
y''-4y'+4y=0
\quad\Longrightarrow\quad
(r-2)^2=0.
\]
Por tanto
\[
\boxed{y=(c_1+c_2x)e^{2x}}.
\]
\end{ejemplo}

\begin{ejemplo}
\[
y'''-3y''+3y'-y=0.
\]
Como el polinomio es $(r-1)^3$,
\[
\boxed{y=e^x(c_1+c_2x+c_3x^2)}.
\]
\end{ejemplo}

\subsection{Raíces complejas}
Si $r=\alpha\pm i\beta$, las soluciones reales son
\[
e^{\alpha x}\cos(\beta x),\qquad
e^{\alpha x}\sin(\beta x).
\]

\begin{ejemplo}
\[
y''-2y'+5y=0.
\]
Las raíces son $1\pm2i$. Entonces
\[
\boxed{y=e^x(c_1\cos2x+c_2\sin2x)}.
\]
\end{ejemplo}

\begin{ejemplo}[PVI]
Resolver
\[
y''-2y'+5y=0,\qquad y(0)=1,\quad y'(0)=0.
\]
De $y(0)=1$ resulta $c_1=1$ y de $y'(0)=0$ resulta $c_2=-1/2$. Así,
\[
\boxed{y=e^x\left(\cos2x-\frac12\sin2x\right)}.
\]
\end{ejemplo}

\section{Ecuaciones no homogéneas}
Para
\[
\mathcal L(y)=g(x)
\]
la estructura es
\[
\boxed{y=y_h+y_p}.
\]

\subsection{Coeficientes indeterminados}
\begin{ejemplo}[Término polinomial]
\[
y''-3y'+2y=x.
\]
La homogénea es $c_1e^x+c_2e^{2x}$. Con $y_p=Ax+B$:
\[
2Ax+2B-3A=x,
\]
de donde $A=1/2$, $B=3/4$. Por tanto
\[
\boxed{y=c_1e^x+c_2e^{2x}+\frac{x}{2}+\frac34}.
\]
\end{ejemplo}

\begin{ejemplo}[Exponencial]
\[
y''-3y'+2y=e^{3x}.
\]
Con $y_p=Ae^{3x}$:
\[
(9-9+2)A=1,
\]
por lo que $A=1/2$. Así,
\[
\boxed{y=c_1e^x+c_2e^{2x}+\frac12e^{3x}}.
\]
\end{ejemplo}

\begin{ejemplo}[Resonancia]
\[
y''-3y'+2y=e^x.
\]
Como $e^x$ ya aparece en $y_h$, probamos $y_p=Axe^x$. Se obtiene $A=-1$, luego
\[
\boxed{y=c_1e^x+c_2e^{2x}-xe^x}.
\]
\end{ejemplo}

\begin{ejemplo}[Forzamiento trigonométrico]
\[
y''+y=\cos2x.
\]
Tomando $y_p=A\cos2x+B\sin2x$:
\[
-3A\cos2x-3B\sin2x=\cos2x.
\]
Así $A=-1/3$, $B=0$, y
\[
\boxed{y=c_1\cos x+c_2\sin x-\frac13\cos2x}.
\]
\end{ejemplo}

\section{Cauchy--Euler}
Para
\[
ax^2y''+bxy'+cy=0,\qquad x>0,
\]
se prueba $y=x^r$. La ecuación indicial es
\[
ar(r-1)+br+c=0.
\]

\begin{ejemplo}
\[
x^2y''-3xy'+4y=0.
\]
La ecuación indicial es $(r-2)^2=0$, así que
\[
\boxed{y=x^2(c_1+c_2\ln x)}.
\]
\end{ejemplo}

\begin{ejemplo}
\[
x^2y''+xy'+y=0.
\]
La ecuación indicial es $r^2+1=0$. Por tanto
\[
\boxed{y=c_1\cos(\ln x)+c_2\sin(\ln x)}.
\]
\end{ejemplo}

\section{Cómo elegir el método}
\begin{enumerate}
\item ¿Es homogénea? Resuelve primero la homogénea.
\item ¿Los coeficientes son constantes? Usa el polinomio característico.
\item ¿Aparecen $x^2y''$, $xy'$ y $y$? Prueba Cauchy--Euler.
\item ¿El término derecho es polinómico, exponencial o trigonométrico? Prueba coeficientes indeterminados.
\item ¿Hay resonancia? Multiplica el ensayo particular por una potencia de $x$.
\item ¿Hay datos iniciales? Determina las constantes al final.
\item ¿Hay datos en dos extremos? Resuelve el sistema de condiciones de frontera.
\end{enumerate}

\begin{center}
\begin{tabular}{>{\raggedright}p{0.30\textwidth}p{0.52\textwidth}}
\toprule
\textbf{Lo que ves} & \textbf{Qué probar primero}\\
\midrule
Coeficientes constantes & ecuación característica $p(r)=0$\\
Raíz $r$ de multiplicidad $m$ & $e^{rx},xe^{rx},\ldots,x^{m-1}e^{rx}$\\
Raíces $\alpha\pm i\beta$ & $e^{\alpha x}\cos\beta x,\ e^{\alpha x}\sin\beta x$\\
$ax^2y''+bxy'+cy$ & sustitución $y=x^r$\\
Polinomio en $x$ & ensayo polinomial\\
$e^{ax}$ & ensayo $Ae^{ax}$\\
$\cos bx,\sin bx$ & ensayo $A\cos bx+B\sin bx$\\
Resonancia & multiplicar el ensayo por $x^s$\\
Datos en un punto & problema de valores iniciales\\
Datos en dos puntos & problema de frontera\\
\bottomrule
\end{tabular}
\end{center}


\section{Un ejemplo completo}

Consideremos
\[
y''-4y'+13y=e^{2x}\cos(3x).
\]

\textbf{Paso 1. Homogénea.}
\[
r^2-4r+13=(r-2)^2+9,
\]
por lo que
\[
y_h=e^{2x}(c_1\cos3x+c_2\sin3x).
\]

\textbf{Paso 2. Particular.}
El término derecho tiene la misma frecuencia y el mismo factor exponencial que
la solución homogénea. Hay resonancia. El ensayo correcto es
\[
y_p=xe^{2x}(A\cos3x+B\sin3x).
\]
Al sustituir se obtiene
\[
A=0,\qquad B=\frac16.
\]
Por tanto
\[
\boxed{
y=e^{2x}(c_1\cos3x+c_2\sin3x)
+\frac{x}{6}e^{2x}\sin3x}.
\]

Este ejemplo resume el procedimiento: resolver la homogénea, identificar el tipo
del término forzante, detectar resonancia, construir $y_p$ y comprobar.

\section{Errores frecuentes}

\begin{enumerate}
\item \textbf{Raíz repetida.} Si $r$ tiene multiplicidad $m$, aparecen
$e^{rx},xe^{rx},\ldots,x^{m-1}e^{rx}$.

\item \textbf{Raíces complejas.} Para $r=\alpha\pm i\beta$, una base real es
\[
e^{\alpha x}\cos\beta x,\qquad e^{\alpha x}\sin\beta x.
\]

\item \textbf{Resonancia.} Si el ensayo particular pertenece a la solución
homogénea, hay que multiplicarlo por una potencia adecuada de $x$.

\item \textbf{Constantes prematuras.} Conviene obtener primero la solución general
y aplicar después las condiciones iniciales o de frontera.

\item \textbf{Sustitución automática.} $y=x^r$ es especialmente natural para
Cauchy--Euler; no es un método universal.

\item \textbf{No comprobar.} Sustituir la solución final en la ecuación original
es una verificación rápida y muy eficaz.
\end{enumerate}

\section{Mini--colección de ejemplos}

\begin{ejemplo}
\[
y''+9y=0
\quad\Longrightarrow\quad
\boxed{y=c_1\cos3x+c_2\sin3x}.
\]
\end{ejemplo}

\begin{ejemplo}
\[
y''-8y'+16y=0
\quad\Longrightarrow\quad
\boxed{y=(c_1+c_2x)e^{4x}}.
\]
\end{ejemplo}

\begin{ejemplo}
\[
y''+4y'+13y=0
\quad\Longrightarrow\quad
\boxed{y=e^{-2x}(c_1\cos3x+c_2\sin3x)}.
\]
\end{ejemplo}

\begin{ejemplo}
\[
y''-y=e^{2x}
\quad\Longrightarrow\quad
\boxed{y=c_1e^x+c_2e^{-x}+\frac13e^{2x}}.
\]
\end{ejemplo}

\begin{ejemplo}
\[
x^2y''+xy'-4y=0
\quad\Longrightarrow\quad
r^2-4=0,
\]
por lo que
\[
\boxed{y=c_1x^2+c_2x^{-2}}.
\]
\end{ejemplo}

\section{Ejercicios}
\begin{enumerate}
\item Resolver $y''-7y'+12y=0$.
\item Resolver $y''+6y'+9y=0$.
\item Resolver $y''+4y=0$ con $y(0)=2$, $y'(0)=-1$.
\item Resolver $y''-2y'+5y=0$ con $y(0)=0$, $y'(0)=1$.
\item Resolver $y''-y=x^2$.
\item Resolver $y''+4y=8\cos2x$.
\item Resolver $y''-4y'+4y=e^{2x}$.
\item Resolver $x^2y''-xy'-3y=0$ para $x>0$.
\item Determinar si $\{1,x,x^2\}$ es linealmente independiente.
\item Determinar si $\{e^x,xe^x,e^{2x}\}$ es linealmente independiente.
\item Hallar el Wronskiano de $\{\cos x,\sin x\}$.
\item Resolver $y''-5y'+6y=0$ con $y(0)=1$, $y(1)=0$.
\item Estudiar el PVF $y''+\pi^2y=0$, $y(0)=0$, $y(1)=0$.
\item Resolver $y'''-6y''+11y'-6y=0$.
\item Para $y''+2y'+y=0$, construir una base de soluciones y verificar su independencia.
\end{enumerate}

\section*{Resumen de fórmulas}

Para una ecuación lineal homogénea con coeficientes constantes,
\[
a_ny^{(n)}+\cdots+a_1y'+a_0y=0,
\]
se busca $y=e^{rx}$ y se resuelve $p(r)=0$.

\[
\begin{array}{c|c}
\text{raíces} & \text{soluciones}\\
\hline
r_1,\ldots,r_n\text{ distintas} & e^{r_1x},\ldots,e^{r_nx}\\
r\text{ de multiplicidad }m & e^{rx},xe^{rx},\ldots,x^{m-1}e^{rx}\\
\alpha\pm i\beta & e^{\alpha x}\cos\beta x,\ e^{\alpha x}\sin\beta x
\end{array}
\]

Para una ecuación no homogénea,
\[
\boxed{y=y_h+y_p}.
\]
La dificultad principal está en elegir correctamente $y_p$ y detectar resonancia.

\[
\boxed{
\text{clasificar}
\longrightarrow
\text{resolver la homogénea}
\longrightarrow
\text{buscar una particular}
\longrightarrow
\text{aplicar condiciones}
\longrightarrow
\text{comprobar}
}
\]
\end{document}
